\chapter*{Preface}
\addcontentsline{toc}{chapter}{Preface}
\label{ch:preface}

\section*{About This Manual}

This manual documents SIRCAS, the SIRC Audio System: the sound processor that plays eight-bit samples through six
voices and sends a stereo signal to a digital-to-analogue converter. SIRCAS is the sound system of the SIRCUS console,
but it is a component in its own right, in a package of its own, with a memory of its own. It is a companion to the
\archmanual{}, which defines the SIRC-1 CPU, to the \hwmanual{}, which defines the CPU package and bus, to the
\asmmanual{}, which defines the language used to program them, and to the \videomanual{}, which defines the display
adapter that shares the console's bus.

SIRCAS is a peripheral on a 16-bit bus, and any 16-bit processor that can supply a chip select, a bus strobe and a
read/write line can drive it. It was designed to work with the SIRC-1, so the signal names, the interrupt arrangements
and the programming examples in this manual are those of the SIRC-1.

It is written for:

\begin{itemize}
	\item Programmers writing music and sound code for SIRCAS
	\item Authors of the tools that prepare samples, instruments and sequences
	\item Engineers implementing the device, in hardware or in software
\end{itemize}

\section*{Assumed Knowledge}

This manual assumes familiarity with sampled-sound hardware: sample playback with a fixed-point position and step,
envelopes, panning and the state-variable filter. Readers who need the interrupt mechanism that delivers the
sequencer's completion signal to the CPU should read the exception chapters of the \archmanual{}.

\section*{Revision Status}

This is a draft of the \audiomanual{}, not for distribution. Anything not yet decided in the manual is marked ``not yet
specified''; the manual does not invent values for it.
